\section{Lean verification and its remaining boundaries}\label{sec:lean}

\subsection{The retained ascending pilot}
The original pilot is a standalone Lean 4.33.0 project using the standard library without Mathlib~\cite{Lean4,LeanReference}. It formalizes the ascending mathematical fragment above. Certificate constants are exported from the saved research certificate as data; handwritten definitions and general lemmas determine what those constants must prove.

\begin{center}\small
\begin{tabularx}{\textwidth}{@{}lX@{}}
\toprule
Module & Checked content\\
\midrule
\texttt{RowRules} & Union/intersection on $\mathcal P(\{0,1\})$, inductive forced-cell rules, soundness, atomic uniqueness, and row-kernel compatibility.\\
\texttt{Gluing} & Simulation and invariant preservation along every finite column list; soundness of a closed certificate.\\
\texttt{Core}, \texttt{Data} & Source-cell phases and their factor; the complete six-column alphabet; row, state, and edge tables.\\
\texttt{Checked} & Atomic agreement, row uniqueness, recovery of the native action from completed rows, closure, and sign conditions.\\
\texttt{WordSemantics} & Numerical words, lower-part bounds, signed comparisons, output-mask membership, and the all-width successor theorem.\\
\texttt{SourceBridge} & Equality of observed numerical states and output traces after any number of modeled nonsign source cells.\\
\texttt{Audit} & Expected axiom dependencies of the principal theorems, enforced at build time.\\
\bottomrule
\end{tabularx}
\end{center}

\begin{theorem}[Lean pilot scope, K10]\label{thm:leanpilot}
The distributed Lean source proves the masked-successor result for every list of legal nonsign columns and each allowed sign completion satisfying its witness premise. It proves the source-word factor for the explicitly modeled nonsign repair iteration. Their names are
\begin{center}\normalfont\small\ttfamily
QKF.mask\_successor\_all\_widths\\
QKF.source\_word\_factor.
\end{center}
The length lemma identifies the complete signed width as the list length plus one.
\end{theorem}
\begin{proof}
The checked certificate supplies the initial state, every transition, and every sign obligation. General gluing yields the invariant for arbitrary lists. A separate relation connects observations to accumulated natural-number values and proves each lower part less than its power-of-two bound. Signed comparison translates the relational assertions into $g<t\le x$, while trace induction establishes mask membership. For the source bridge, local phase factorization commutes with each numerical step and is lifted by the simulation theorem. These are the dependencies of the named Lean proofs.
\end{proof}

Finite obligations use \texttt{decide}; elementary integer bridges use \texttt{omega}. The audit reports no axioms for forced-rule soundness and atomic row uniqueness. The main numerical theorem uses only the standard \texttt{propext} and \texttt{Quot.sound}; the source-word factor uses \texttt{propext}. No additional axioms, proof placeholders, \texttt{native\_decide}, or \texttt{bv\_decide} are used. Exact expected axiom lists are enforced by \texttt{\#guard\_msgs}.

The pilot does not formalize both foundational papers in full, the 49-state descending certificate, the outer lower-order argument, the complete Java helper, or the Python checker. Its source bridge concerns the stated operational cell model; translating actual Java and host-long arithmetic remains a separate obligation. Lean verification improves confidence in this fragment without changing whole-program coverage counts.

The package has been built from a fresh extracted copy, and three numerical examples are checked. Four isolated negative controls change an initial state, an edge, a row cell, and a proof; all are rejected at the corresponding certificate or axiom-audit gate. These controls exercise the build and trust checks; proof terms establish the general results.


\subsection{Numeric targets and conditional caller rules}
A separate pinned project, \artifact{research/lean_targets/MASKED_WORDS_RU.md}, proves an exact numeric cyclic-successor theorem from an accepted target certificate. Every legal competitor is represented by a bit list of the same length; the emitted result depends only on the input projection, not on the chosen competitor. The theorem \texttt{exact\_cyclic\_successor} states membership, the least strictly greater legal value when one exists, and return to the minimum otherwise, for every positive input length. Its two retained source factors have two source classes, nine joint states, and 54 transitions each. These figures describe that target interface, not a replacement of the eight-state certificate in the original pilot.

\artifact{research/lean_targets/CEILING_COMPOSITION_RU.md} formalizes the next composition. An exact floor below a bound, followed by an exact cyclic successor from precisely that floor, gives the least admissible value at or above the bound when the caller supplies the required extrema. The empty result is an explicit option distinct from numeric zero. The floor contract is a theorem hypothesis; the proof does not validate the descending Java region by assuming a saved Python verdict. Mask equality and reseeding preserve the legal numeric carrier. Caller lemmas in \artifact{research/lean_targets/CREATE_CALLER_RU.md} establish sign-bucket order, common-prefix preservation, extrema monotonicity, nonwrapping witnesses, and the conditional final-pass argument. Java parsing and the establishment of the source-level premises remain separate.

\subsection{F1a: the decoded typed positive ground calculus}
The standalone project \artifact{formal/ground/README_RU.md} supplies a computable Lean checker for positive chronological ground-DAG certificates. Lean checks backward node references, symbol existence, arity, sorts, equations, proof events, argument paths, all positive goals, and recomputed depths. Every congruence premise refers to the accepted earlier prefix. Unused events are also checked.

\begin{theorem}[Formal ground soundness]\label{thm:formalground}
For an accepted decoded request and certificate, every accepted goal has equal endpoint values in every many-sorted interpretation satisfying its declared ground equations $E$. The two endpoints have the same sort and actual values in that carrier.
\end{theorem}
\begin{proof}
\texttt{QKFGround.checkRaw\_sound} combines DAG admission with the positive proof checker. \texttt{check\_sound} proves preservation of equality for paths, congruence arguments, chronological events, and goals in a total algebra. \texttt{typed\_check\_sound} specializes this to the many-sorted interpretation, and \texttt{eval\_typed} shows that admitted terms evaluate to real values of the declared sort. \texttt{accepted\_goal\_values} packages the resulting endpoint guarantee. These are checked general Lean proofs, not deductions from a finite set of exported examples.
\end{proof}

Empty sorts are permitted when the signature and interpretation allow them; no artificial default value is introduced for a sort. The theorem therefore does not mistake two failed evaluations for equal semantic values. Checked depth establishes validity at the reported horizon, not minimum depth. Negative models, exact-threshold minimality, search completeness, and source-rule truth are outside F1a. Its input is a Lean \texttt{Request}, not unparsed JSON bytes; the Python-to-Lean name mapping and exporter remain an adapter boundary.

\subsection{F1b: chronological import and conservative extension}
The project \artifact{formal/composition/README_RU.md} uses the preceding ground kernel unchanged. Its decoded packet contains a common signature, explicit native premises $E$, chronological entries, and independently specified consumer endpoints. A native entry refers to an actual member of $E$. A derived entry supplies a ground certificate whose axioms exactly match selected earlier equalities, in their declared order. A \texttt{via} reference requires a previously established $a=b$ and an exact residual proof $b=c$ to produce $a=c$. Consumers point to accepted equalities and repeat their exact endpoints.

\begin{theorem}[Formal checked composition]\label{thm:formalcomposition}
If the composition checker accepts, all imported equalities $L$, including imports unused by the final consumer, and all accepted consumer equalities hold in every interpretation satisfying the original $E$. In particular,
\[
 \operatorname{Models}(E\cup L)=\operatorname{Models}(E).
\]
For typed consumers, both endpoints have the same actual semantic value.
\end{theorem}
\begin{proof}
Induction over chronological entries uses ground soundness for each selected earlier basis. Explicit residual composition uses transitivity. Every model of $E$ therefore satisfies the imported equalities, while every model of $E\cup L$ satisfies $E$. Exact consumer endpoints yield their stated equalities. The Lean results are \texttt{all\_admitted\_sound}, \texttt{conservative\_extension}, \texttt{check\_sound}, and \texttt{accepted\_consumer\_values} in the \texttt{QKFComposition} namespace.
\end{proof}

Together F1a and F1b formalize the decoded typed positive proof calculus and checked-premise composition. Native facts remain explicit assumptions, carrying evidence and scope checked by Python. Source epoch, guard, width, cached-request identity, JSON decoding, frontend semantics, and interface-to-runtime correspondence are not established by these Lean theorems. Exact hashes bind the retained inputs without proving that remaining semantic correspondence.

The F1b validation retains 12 positive packets with 56 consumer obligations, including four original SDK packets from the ordinary/query and no-lemmas/reuse matrix. It audits 12 general theorems and 29 named example results and rejects 17 corrupted Lean packets. The accepted formal projects use Lean 4.33.0 and Std, with only the audited standard \texttt{propext} and \texttt{Quot.sound} dependencies where needed; proof-hole and extra-axiom controls test that boundary. Exact results and scope are retained in \href{https://github.com/Len989/QKF-Certifier/blob/ab43bbe11ed99ff5e21808a34351c37915a1ec2f/formal/composition/REPORT_RU.md}{the F1b validation report} and \href{https://github.com/Len989/QKF-Certifier/blob/e210e63dbf19f045f3ad84b528a22810c1fb4d2e/formal/ground/REPORT_RU.md}{the F1a validation report}, at their separate results-only commits. None of these counts enlarges the external program corpus.
